%% LyX 2.4.4 created this file.  For more info, see https://www.lyx.org/.
%% Do not edit unless you really know what you are doing.
\documentclass[english]{article}
\usepackage[T1]{fontenc}
\usepackage[latin9]{inputenc}
\usepackage{enumitem}
\usepackage{amsmath}
\usepackage{amsthm}
\usepackage{cancel}
\usepackage{geometry}
\geometry{verbose,tmargin=1in,bmargin=1in,lmargin=1in,rmargin=1in}

\makeatletter
%%%%%%%%%%%%%%%%%%%%%%%%%%%%%% Textclass specific LaTeX commands.
\numberwithin{equation}{section}
\numberwithin{figure}{section}
\newlength{\lyxlabelwidth}      % auxiliary length 

%%%%%%%%%%%%%%%%%%%%%%%%%%%%%% User specified LaTeX commands.
\usepackage{hyperref}
\usepackage[usenames,dvipsnames,table]{xcolor} %adds additional color options
\hypersetup{colorlinks=true, urlcolor=black,linkcolor=black}

\usepackage{fancyhdr}
\usepackage{lastpage}
\pagestyle{fancy}
\usepackage{enumitem}
\usepackage{ifthen}
%sets math fonts to be more readable
%\everymath=\expandafter{\the\everymath\displaystyle}
%\DeclareMathSizes{10}{10}{10}{10}
\usepackage{multicol}

\usepackage{tikz}
\usetikzlibrary{plotmarks}
\usepackage{pgfplots}
\pgfplotsset{compat=newest}

\usepackage{calc}
\usepackage[nomessages]{fp}

\usepackage{advdate}    % Advancing/saving dates
\usepackage{datetime}   % Dates formatting
\usepackage{datenumber} % Counters for dates

\usepackage{fancybox} %%ovalbox and fancy box settings

\usepackage[super]{nth} %easy 1st, 2nd, etc

\usepackage{forest} %%for tree diagram
% Added by lyx2lyx
\setlength{\parskip}{\smallskipamount}
\setlength{\parindent}{0pt}

\makeatother

\usepackage{babel}
\begin{document}
\renewcommand{\author}{Dr.\,James Ashe}

\newcommand{\classNum}{107}

\newcommand{\classSec}{7 and 10}

\newcommand{\nameType}{Solutions}

\newcommand{\examType}{Quiz 4.4}

\renewcommand{\date}{\formatdate{12}{11}{2014}}

%%First page's header%%%%%%%%%%%%%%%%%%%%%%
\fancypagestyle{firststyle} %%header settings for first pagr
{
	\fancyhead[L]{MTH \classNum{}(\classSec{}) \author{}\\
	\textbf{\examType{}} \date{}}
	\fancyhead[C]{\textbf{\nameType}}
	\setlength{\headsep}{14pt}
}
%%%%%%%%%%%%%%%%%%%%%%%%%%%%%%%%%%%%%%%%%%

%%Next page's header%%%%%%%%%%%%%%%%%%%%%%%
%\setlist{nolistsep}
\lhead{MTH \classNum{}(\classSec{})} 
\chead{\textbf{\nameType}} 
\cfoot{\thepage{}~of~\pageref{LastPage}} 
\setlength{\headsep}{5pt} 
%%%%%%%%%%%%%%%%%%%%%%%%%%%%%%%%%%%%%%%%%%%

%%Various settings; see the preamble for other settings%%%%%
%%\setenumerate[0]{leftmargin=12pt,itemindent=0pt,labelwidth=0pt}
%\setlist[2]{noitemsep}
%\setlist{nosep}
\setlist{itemsep=3pt,topsep=0pt}
\renewcommand{\labelenumi}{\textbf{\arabic{enumi}.}}
\renewcommand{\labelenumii}{\textbf{\alph{enumii}.}}
\thispagestyle{firststyle}
%%%%%%%%%%%%%%%%%%%%%%%%%%%%%%%%%%%%%%%%%%
\newcommand{\xmax}{10}
\newcommand{\xmin}{-10}
\newcommand{\xstep}{0} %must be xmin+n
\newcommand{\ymax}{10}
\newcommand{\ymin}{-10}
\newcommand{\ystep}{0} %must be ymin+n
\newcommand{\xspace}{\hspace*{40pt}}
\newcommand{\yspace}{\vspace*{1.4cm}}
\newcommand{\rowColor}{gray!35}
\large

\newcommand{\gauss}[2]{1/(#2*sqrt(2*pi))*exp(-((x-#1)^2)/(2*#2^2))}

\newcommand{\normalcdf}[4]{ %mean,std,z-left,z-right
\begin{tikzpicture}[scale=1.25]

\pgfmathsetmacro{\mn}{#1}
\pgfmathsetmacro{\sd}{#2}
\pgfmathsetmacro{\za}{#1+#2*1}
\pgfmathsetmacro{\zb}{#1+#2*2}
\pgfmathsetmacro{\zc}{#1+#2*3}

\pgfmathsetmacro{\zna}{#1-#2*1}
\pgfmathsetmacro{\znb}{#1-#2*2}
\pgfmathsetmacro{\znc}{#1-#2*3}

\pgfmathsetmacro{\zleft}{#3}
\pgfmathsetmacro{\zright}{#4}
\pgfmathsetmacro{\zsL}{(#3-#1)/#2} %z-score
\pgfmathsetmacro{\zsR}{(#4-#1)/#2} %z-score
\pgfmathsetmacro{\xCenter}{\zsL/2.5+\zsR/2.5} %center using z-score
\pgfmathsetmacro{\yCenter}{((1/(1*sqrt(2*pi))*exp(-((\xCenter-0)^2)/(2*1^2)))/2}
\pgfkeys{/pgf/fpu} %for large numbers
\pgfmathparse{100*(1/(1 + exp(-0.07056*((\zsR-0)/1)^3 - 1.5976*(\zsR-0)/1))-1/(1 + exp(-0.07056*((\zsL-0)/1)^3 - 1.5976*(\zsL-0)/1)))} %area from z-left to z-right
\edef\area{\pgfmathresult} 
\pgfkeys{/pgf/fpu=false}

\scriptsize
\begin{axis}[ 
	no markers, 
	domain=0:10, 
	%samples=100, 
	axis lines*=left, 
	%xlabel=Standard deviations, 
	%ylabel=Frequency, 
	height=2.9cm, %change/remove this scaling to get 1:1 pic
	width=6cm, 
	hide y axis,
	%ticks=none,
	xtick={-3,-2,-1,0,1,2,3},
	xticklabels={\pgfmathprintnumber[precision=0]{\znc},\pgfmathprintnumber[precision=2]{\znb},\pgfmathprintnumber[precision=2]{\zna},#1,\pgfmathprintnumber[precision=0]{\za},\pgfmathprintnumber[precision=2]{\zb},\pgfmathprintnumber[precision=2]{\zc}},
	%xticklabels={-25, -15, -5, 5, 15, 25, 35},
	%ytick=empty, 
	enlargelimits=false, 
	clip=false, 
	axis on top, 
	%grid = major
	]

	\addplot [fill=blue!50, draw=red,  domain=\zsL:\zsR] {\gauss{0}{1}} \closedcycle;
		\node (arrowEnd) at (axis cs: \xCenter, \yCenter) {};
	\addplot [draw, thick, smooth, domain=-3:3] {\gauss{0}{1}} \closedcycle;
		\node[anchor=south east] (arrowStart) at (current bounding box.east) {\pgfmathprintnumber[fixed,precision=1]{\area}\%};
	\draw[->,thick,black] (arrowStart) to (arrowEnd);
	
	%\addplot [fill=orange!20, draw=none, domain=2:3] {\gauss{0}{1}} \closedcycle; 
	%\addplot [fill=blue!20, draw=none, domain=-2:-1] {\gauss{0}{1}} \closedcycle; 
	%\addplot [fill=blue!20, draw=none, domain=1:2] {\gauss{0}{1}} \closedcycle; 
	%\only<1->{\node[above] at (axis cs: -0.5, 0.025) {$34.1\%$};} 
	%\only<1->{\node[above] at (axis cs: 0.5, 0.025) {$34.1\%$};} 
	%\only<1->{\node[above] at (axis cs: -1.5, 0.025) {$13.6\%$};}
	%\only<1->{\node[above] at (axis cs: 1.5, 0.025) {$13.6\%$};}
	%\only<1->{\node[above] at (axis cs: 2.5, 0.05) {$2.1\%$};}
	%\only<1->{\node[above] at (axis cs: -2.5, 0.05) {$2.1\%$};}
	%\only<1->{ \addplot[<->,black,very thick] coordinates {(-1,0.4) (1,0.4)};} 
	%\only<1->{\node[above] at (axis cs: 0, 0.4) {$68.2\%$};} 
	%\only<1->{\addplot[<->,black,very thick] coordinates {(-2,0.3) (2,0.3)};}
	%\only<1->{\node[above] at (axis cs: 0, 0.3) {$95.4\%$};}
	%\only<1->{\addplot[<->,black, very thick] coordinates {(-3,0.2) (3,0.2)};} 
	%\only<1->{\node[above] at (axis cs: 0, 0.2) {$99.7\%$};}
	\end{axis} 
\end{tikzpicture}
}

\emph{Unless stated otherwise, all problems are worth equal value
and answers should be stated as decimals rounded to three decimal
places. You must show your work to receive credit. }
\begin{enumerate}[resume]
\item A population is normally distributed with mean 24.7 and standard deviation
2.3.

\begin{enumerate}
\item Draw a sketch of the bell curve labeling the mean and the intervals
representing 1, 2, and 3 standard deviations of the mean.

\begin{enumerate}
\item []\normalcdf{24.7}{2.3}{18}{31.6}
\end{enumerate}
\item Find the percentage of data lies in each of the intervals.
\end{enumerate}
\end{enumerate}
\begin{multicols}{3}

\scalebox{0.9}{\normalcdf{24.7}{2.3}{22.4}{27}}

\scalebox{0.9}{\normalcdf{24.7}{2.3}{20.1}{29.3}}

\scalebox{0.9}{\normalcdf{24.7}{2.3}{18}{33.9}}

\end{multicols}
\begin{enumerate}[resume]
\item Find the following probabilities, 

\begin{enumerate}
\item $p(z<1.75)$\begin{multicols}{2}

\begin{enumerate}
\item []$p(z<1.75)\approx95.99\%$ \columnbreak
\item []\normalcdf{0}{1}{-4}{1.75}
\end{enumerate}
\end{multicols}
\item $p(0<z<1.75)$\begin{multicols}{2}

\begin{enumerate}
\item []$p(z<1.75)-p(z<0)$
\item []$\approx95.99\%-50\%\approx45.99\%$ \columnbreak
\item []\normalcdf{0}{1}{0}{1.75}
\end{enumerate}
\end{multicols}
\item $p(-.5<z<1.75)$\begin{multicols}{2}

\begin{enumerate}
\item []$p(z<1.75)-p(z<-.5)$ 
\item []$\approx95.99\%-30.85\%\approx65.14\%$ \columnbreak
\item []\normalcdf{0}{1}{-.5}{1.75}
\end{enumerate}
\end{multicols}
\item $p(-1.5<z<-.5)$\begin{multicols}{2}

\begin{enumerate}
\item []$p(z<-.5)-p(z<-1.5)$ 
\item []$\approx30.85\%-6.68\%\approx24.17\%$ \columnbreak
\item []\normalcdf{0}{1}{-1.5}{-.5}
\end{enumerate}
\end{multicols}
\end{enumerate}
\end{enumerate}
\newpage
\begin{enumerate}[resume]
\item A population is normally distributed with a mean 250 and standard
deviation 24. Find the following probabilities. 

\begin{enumerate}
\item $p(x<200)$\begin{multicols}{2}

\begin{enumerate}
\item []$z=\frac{200-250}{24}=2.08\bar{3}$
\item []$p(z<-2.08)\approx1.88\%$\columnbreak
\item []\normalcdf{250}{24}{154}{200}
\end{enumerate}
\end{multicols}\vfill
\item $p(200<x<250)$\begin{multicols}{2}

\begin{enumerate}
\item []$z_{200}=\frac{200-250}{24}=2.08\bar{3}$ and
\item []$z_{250}=\frac{250-250}{24}=0$
\item []$p(-2.08<z<0)$\\
$\approx50\%-1.88\%\approx48.12\%$\columnbreak
\item []\normalcdf{250}{24}{200}{250}
\end{enumerate}
\end{multicols}\vfill
\item $p(200<x<300)$\begin{multicols}{2}

\begin{enumerate}
\item []$z_{200}=\frac{200-250}{24}=-2.08\bar{3}$ and
\item []$z_{300}=\frac{300-250}{24}=2.08\bar{3}$
\item []$p(-2.08<z<2.08)$\\
$\approx98.12\%-1.88\%\approx96.24\%$\columnbreak
\item []\normalcdf{250}{24}{200}{300}
\end{enumerate}
\end{multicols}\vfill
\item $p(154<x<200)$\begin{multicols}{2}

\begin{enumerate}
\item []$z_{200}=\frac{200-250}{24}=-2.08\bar{3}$ and
\item []$z_{154}=\frac{154-250}{24}=-4$
\item []$p(-4<z<-2.08)$\\
$\approx1.88\%-0\%\approx1.88\%$\columnbreak
\item []\normalcdf{250}{24}{154}{200}
\end{enumerate}
\end{multicols}\vfill
\end{enumerate}
\end{enumerate}

\end{document}
